\subsection*{Mercy answers by giving life into communion\enspace\properrefs{Int., Gosp., Off., Comm.}}

The Introit's servant raises his soul and asks for an answering ear. The
Gospel discloses the divine initiative toward a need that has not found
words: Christ sees the widow, is moved with mercy, and commands her dead son
to rise. At the Offertory, waiting becomes hearing and a new canticle; at
Communion, Christ's own flesh is given for the world's life. The movement is
not merely from petition to result but from isolated distress into received
and communicated life.

Augustine hears Psalm~85 as the prayer of the needy whole Christ and
Psalm~39's new song as the song of renewed humanity. Cyril reads both the
touch at Nain and the flesh in John~6 through the incarnate Word's
life-giving power. Their reasoning joins the ritual moments while the
psalmist, the widow, and the restored youth each retain a distinct voice
before the one Lord whose mercy hears, whose word raises, and whose gift
communicates life.

Ambrose carries the Gospel into the Church: the widow figures Mother Church,
her tears intercede for a son, the bier figures the cross, and restored
speech becomes confession (\emph{Expositio in Lucam} V.89--92). Augustine
reads the one bread of John~6 as ecclesial incorporation in charity
(\emph{Tractate on John} 26 §§12--13). The formulary therefore makes renewed
life public and ecclesial. Favorable temporal outcomes are no simple measure
of whether God has heard prayer.

\subsection*{The gift that precedes also commands\enspace\properrefs{Coll., Ep., Sec., Postcomm.}}

The Collect confesses that the Church cannot stand safely without continuing
mercy. Yet Galatians speaks through imperatives: walk, restore, bear, examine,
share, sow, persevere, do good. The Secret asks sacramental protection and the
Postcommunion asks that the gift's operation possess mind and body and
precede self-directed judgment. Divine action is first, intimate, and
sustaining; human action remains responsible, patient, and concrete.

Chrysostom locates Spirit-led freedom in loving service rather than in
lawless autonomy. Aquinas makes restoration of the fallen an exercise of
patience, aid, prayer, and charity. Augustine distinguishes mutual burden
bearing now from the personal account at judgment. Thus Challoner's ``Bear ye
one another's burdens'' does not dissolve the later assignment of each
person's own burden. The law
of Christ gives mutual care its form, and final accountability gives it
seriousness.

The Secret adds an enemy without excusing the agent. Diabolical assaults are
real, yet the Epistle also names vainglory, envy, self-deception, corruptible
sowing, and weariness. The faithful ask to be guarded so that they may
persevere, not so that sacramental action may replace repentance, prudence,
or charity. The heavenly gift possesses the whole person by healing and
directing agency, not by destroying it.

\subsection*{One formulary coordinates four registers of life\enspace\properrefs{Ep., Gosp., Comm., Postcomm.}}

Paul's moral life in the Spirit bears an eschatological harvest of
everlasting life. Luke's young man returns bodily to mortal life. John~6
speaks of the flesh given for the world's life within promises of abiding and
resurrection on the last day. The Postcommunion places mind and body together
under the heavenly gift. The four elements coordinate conduct,
resuscitation, sacramental participation, and final hope while preserving
their different objects.

Cyril's Nain exegesis treats the miracle as a pledge of the general
resurrection, not the final event itself. His John commentary grounds
life-giving flesh in the personal union of flesh and Word. Augustine
emphasizes inward participation and one-bread incorporation; Chrysostom holds
together faith and bodily Mystery; Aquinas links real flesh, Passion,
restored life, and ecclesial unity. The sacramental reading is therefore
christological and ecclesial rather than a cipher imposed upon Paul's sowing
or Luke's funeral procession.

This distinction also guards the body. Paul's ``flesh'' names a morally
disordered field of sowing, not bodily substance as evil. John names Christ's
flesh as gift, and the Postcommunion asks for the body itself to be possessed
by grace. The liturgy opposes corruption without opposing creation.

\subsection*{Praise widens from one servant into a universal horizon\enspace\properrefs{Int., Coll., Ep., Grad., All., Gosp., Off., Comm.}}

The Introit begins with a singular servant; the Collect names the Church; the
Epistle orders mutual life and good done to all; the Gospel joins two crowds;
the Alleluia acclaims the great King with \latin{super omnem terram}; the
Communion names the world. Praise grows because mercy received cannot remain private. At Nain,
fear becomes glorification and the recognition of divine visitation. In
Psalm~39, a new song appears in the mouth and, in the unappointed continuation,
becomes visible to many.

The Gradual gives this praise a temporal span. Augustine interprets morning
and night through prosperity and adversity; Theodoret reads unceasing praise
of mercy and truthful judgment within Sabbath rest. Their accents show why
praise is neither optimism nor denial of darkness. It announces mercy in the
morning and truth through the night because both belong to providence.

The Alleluia's textual difference yields two distinct patristic horizons.
Augustine comments upon the Clementine \latin{super omnes deos}; Theodoret's
lemma carries \latin{super omnem terram}, the form the Missal appoints.
Covenant, the rejection of idols, and universal dominion all widen the
horizon, while Paul's special care for the household of faith preserves an
ordered ecclesial proximity within good done to all.

\subsection*{Sacramental protection serves transformation\enspace\properrefs{Coll., Ep., Sec., Comm., Postcomm.}}

The Collect asks cleansing, strengthening, and governance; the Secret asks
protection; the Communion gives Christ's flesh for life; the Postcommunion
asks the heavenly gift to act upon mind and body. Defense and transformation
are therefore two aspects of the same sacramental action. The Church is
guarded from what deforms her so that she may be conformed to the life she
receives.

The Epistle describes the visible fruit of that action. Gentleness governs
correction, humility tests the self, charity carries another, responsibility
bears its own load, generosity shares goods, and hope perseveres toward the
harvest. These are not substitutes for sacrament, and the sacrament is not a
substitute for them. Augustine's inward participation, Chrysostom's loving
service, and Aquinas's application of the Passion place moral and
sacramental fruit within one life of grace.

The Trinity Preface supplies the Sunday's normal Trinitarian confession:
Galatians names the Spirit, John~6's context names the Father and Son, and the
Preface worships the Trinity. The life that guards and transforms comes from
the triune God and returns as ecclesial praise.
